\chapter{The Structured-Products Business}\label{ch:dv:the-structured-products-business}

A private-bank client buys a two-year note on an index. It returns all of his capital at maturity, plus 70\% of the
index's rise. The option that pays the 70\% is an ordinary at-the-money call, and at the index's volatility and the
current rate it costs about 12.5 per 100. What decided whether 70 was possible was not that option. It was the price at
which the bank that issued the note borrows. A note is a bond of its issuer, and an issuer that borrows at a high spread
over the risk-free rate discounts the guaranteed 100 more steeply, which leaves more of the client's money to buy options.
With no spread the note could have offered 35\%. At a spread of 2.4\% a year it can offer 70\%. The client was paid for
lending to the bank, and may not have known it. This chapter takes the business apart: its products and wrappers, the
decomposition that prices every note, distribution and its margins and rules, the issuer's funding, and the systematic
indices that products are increasingly written on.

\section{Products and wrappers}

\begin{definition}[Structured product]\label{def:dv:the-structured-products-business:product}
A \emph{structured product}\index{structured product} is an investment whose payoff is a pre-defined function of one or more
underlyings, built by combining a fixed-income instrument with derivatives and sold as a single package.
\end{definition}

\begin{definition}[Structured note]\label{def:dv:the-structured-products-business:note}
A \emph{structured note}\index{structured note} is a \omterm{def:dv:the-structured-products-business:product}{structured product} issued as a debt security of a bank or its issuing
vehicle: the investor holds the issuer's credit risk as well as the payoff's market risk.
\end{definition}

The same payoff can be sold in several wrappers. A note is the most common. A certificate or warrant listed on an exchange
is still a debt obligation of its issuer. The payoff can also sit in a fund, with collateral held separately, or in a
deposit, often with guarantees on part of the capital. It can be written as an over-the-counter swap for institutions, whose
counterparty risk is collateralised. It can even be packaged as an exchange-traded fund (One Quant Book 1, chapter 14) or a
total return swap (Book 1, chapter 6) on a strategy index. The wrapper decides who bears the issuer's credit risk, how the
product is taxed and regulated, and whether a secondary market exists. The payoff is decomposed the same way in all of them.

The two archetypes are the \omterm{def:dv:the-structured-products-business:protected}{capital-protected note} and the yield enhancer.

\begin{definition}[Capital-protected note]\label{def:dv:the-structured-products-business:protected}
A \emph{capital-protected note}\index{capital-protected note} repays at least a stated share of its notional at maturity,
such as 90\% or 100\%, whatever the underlying does, plus a share of the underlying's rise; the protection is only as good as
the issuer's credit.
\end{definition}

\begin{definition}[Reverse convertible]\label{def:dv:the-structured-products-business:rc}
A \emph{reverse convertible}\index{reverse convertible} pays a coupon above the issuer's funding rate and repays par at
maturity unless the underlying has fallen below a strike, in which case it repays the underlying's value (or shares): the
investor has sold a put.
\end{definition}

\begin{definition}[Participation rate]\label{def:dv:the-structured-products-business:participation}
The \emph{participation rate}\index{participation rate} of a note is the share of the underlying's performance it pays above
the protected amount, $p$ in $100\,(\text{protection})+100\,p\,(S_T/S_0-1)^+$.
\end{definition}

Between the archetypes sit the \omterm{def:dv:autocallables:autocallable}{autocallables} of chapter 18, which trade capital protection for coupons, and a long tail of
variations. There are caps and floors on the participation, averaging of the final level (chapter 16), worst-ofs on baskets
(chapter 17), and knock-in barriers (chapter 15). Every variation is a change to the option leg, and each has a price that
the decomposition makes visible.

\section{Anatomy of a note: a bond plus options}

Every note is priced by the same identity. The issue price, 100, equals the issuer's zero-coupon bond for the guaranteed
amounts, plus the options, plus the margin:
\[
100=\underbrace{\pi\cdot100\,e^{-(r+s)T}}_{\text{bond at the issuer's rate}}+\underbrace{p\cdot C}_{\text{options}}+
\underbrace{m}_{\text{margin}},
\]
with $\pi$ the protected share, $r$ the risk-free rate, $s$ the issuer's funding spread (Book 1, chapter 17), $C$ the price of
the option leg per unit of participation, and $m$ the \omterm{def:dv:the-structured-products-business:margin}{structuring margin}. Given any three of $p$, $s$, $m$ and the option's
terms, the identity yields the fourth.

\begin{definition}[Structuring margin]\label{def:dv:the-structured-products-business:margin}
The \emph{structuring margin}\index{structuring margin} is the difference between a note's issue price and the value of its
components at the issuer's own funding curve and the desk's option prices: the revenue shared between the issuer and the
distributors.
\end{definition}

\begin{example}[Seventy percent participation]\label{ex:dv:the-structured-products-business:participation}
A two-year note, 100\% protected, on an index with a 3\% rate and a 1.5\% dividend yield, at chapter 9's two-year
at-the-money volatility of 20.5\%. The at-the-money call costs 12.49 per 100, and the margin is 1.5. At a zero funding
spread the bond costs 94.18, leaving 4.32 for options: a participation of 34.6\%. Each half point of spread adds about 7.5
points of participation. At 1\% it is 49.6\%, at 2\% it is 64.2\%, and at 2.40\% the bond costs 89.76 and the participation reaches
70\% (\cref{fig:dv:the-structured-products-business:participation}).
\end{example}

\begin{omfigure}
\begin{tikzpicture}
\begin{axis}[width=10.5cm,height=5cm,
  xlabel={issuer funding spread (basis points a year)},ylabel={participation (\%)},
  tick label style={font=\small},label style={font=\small},
  xmin=0,xmax=300,ymin=30,ymax=82,grid=major,grid style={gray!25},table/col sep=comma]
\addplot[omThm,very thick,mark=*,mark size=1.4pt] table[x=spread,y=participation]{figdata/derivatives/19-the-structured-products-business/participation.csv};
\draw[gray,dashed] (axis cs:0,70) -- (axis cs:240.3,70) -- (axis cs:240.3,30);
\node[font=\footnotesize,anchor=south east] at (axis cs:235,70.5) {70\% at 240 bp};
\end{axis}
\end{tikzpicture}

\omcaption{The participation a two-year 100\%-protected note can offer, after a 1.5\% margin, against the issuer's funding
spread (rate 3\%, dividend yield 1.5\%, chapter 9's two-year volatility). The spread, not the option, sets the participation:
without it the note could offer 35\%. Data: the tutorial.}\label{fig:dv:the-structured-products-business:participation}
\end{omfigure}

The spread is paid by the investor. A higher spread means a weaker issuer, whose default would take the protection with it. The
70\% is compensation for that risk, and a note from a stronger issuer offering 50\% is not worse value. The same identity prices
the other shapes. With a 1\% spread the budget buys 100\% participation up to a cap of 113.3\%, a two-year return capped at 13.3\%,
instead of 49.6\% uncapped (\cref{fig:dv:the-structured-products-business:payoffs}). With no spread, the cap falls to 108.7\%. The
\omterm{def:dv:the-structured-products-business:rc}{reverse convertible} reverses the flow: the investor sells a put, and the premium is added to the coupon.

\begin{example}[A one-year reverse convertible]\label{ex:dv:the-structured-products-business:rc}
Strike 90\%, one year, the issuer's spread 1\%, the same market. The investor sells $100/0.9$ puts struck at 90\%, worth 4.26, and the
issuer's funding benefit is worth 3.92. After a 1.5 margin the note pays a coupon of 6.96\%, against a risk-free rate of 3\%. If the
index ends below 90\%, the investor receives $100\,S_T/(0.9\,S_0)$ instead of par, and the coupon.
\end{example}

\begin{omfigure}
\begin{tikzpicture}
\begin{axis}[width=10.5cm,height=5cm,
  xlabel={index at maturity (\% of the start)},ylabel={repayment (per 100)},
  tick label style={font=\small},label style={font=\small},
  xmin=50,xmax=160,ymin=90,ymax=145,grid=major,grid style={gray!25},
  legend columns=2,legend cell align=left,legend style={font=\footnotesize,at={(0.5,-0.3)},anchor=north,draw=none,fill=none,/tikz/every even column/.append style={column sep=8pt}},
  table/col sep=comma]
\addplot[omThm,very thick] table[x=perf,y=protected]{figdata/derivatives/19-the-structured-products-business/payoffs.csv};
\addplot[omDef,very thick,dashed] table[x=perf,y=capped]{figdata/derivatives/19-the-structured-products-business/payoffs.csv};
\legend{70\% participation (2.4\% spread),100\% up to a cap (1\% spread)}
\end{axis}
\end{tikzpicture}

\omcaption{Two two-year 100\%-protected notes that each cost the client 100: 70\% participation without a cap, which needs a
2.4\% funding spread, and full participation capped at 113.3\%, which a 1\% spread pays for. Data: the tutorial.}\label{fig:dv:the-structured-products-business:payoffs}
\end{omfigure}

\section{Distribution, margins and regulation}

A note passes through several hands before it reaches its buyer. The structurer designs it, and the trading desk prices and hedges
it. The issuer, often a separate legal entity, provides the balance sheet. The distributor, a private bank, a retail bank or an
independent adviser, sells it. The margin in the identity above is shared among them. The desk keeps a reserve for hedging costs, and
the issuer and the distributor are paid through the note's price.

Regulators focus on what the retail buyer can see. The European Union's regulation on packaged retail investment products requires a
short standardised key information document for every such product. The United States self-regulator FINRA has issued guidance on
products with ``novel, complicated or intricate derivative-like features'', \omterm{def:dv:the-structured-products-business:note}{structured notes} among them.

\begin{dated}{2026-09}{Disclosure rules for retail structured products}\label{dat:dv:the-structured-products-business:rules}
\emph{European Union.} Regulation (EU) No 1286/2014 on packaged retail and insurance-based investment products (PRIIPs) requires
the manufacturer, before a product is made available to retail investors, to draw up a key information document and publish it on
its website. The document is pre-contractual information, must be accurate, fair, clear and not misleading, stands alone, separate
from marketing material, and runs to at most three sides of A4 paper. \emph{United States.} FINRA Regulatory Notice 12-03 (January
2012), ``Heightened Supervision of Complex Products'', asks broker-dealers to apply heightened supervision to products such as
\omterm{def:dv:the-structured-products-business:note}{structured notes}, and follows earlier notices on \omterm{def:dv:the-structured-products-business:product}{structured products}, principal-protected notes and \omterm{def:dv:the-structured-products-business:rc}{reverse convertibles}. Both texts
have been amended or supplemented since; check the current versions.
\end{dated}

For a quant the regulatory text matters in three places. The disclosure's performance scenarios and cost figures are computed from
models, and the model must be the one the desk uses. The margin must be defensible as fair value. And the product's value after issue,
the price at which the issuer buys notes back, must follow from the same pricing as the issue price, less a bid-offer that the
investor was told about.

\section{Issuer funding and the product lifecycle}

A \omterm{def:dv:the-structured-products-business:note}{structured note} is funding for its issuer. The bank receives 100 today and owes the payoff later. If it hedges the payoff with
derivatives, what remains is a loan at its own funding spread. That makes notes one of a bank's funding sources, and it explains why
issuers compete on the rate at which they value their own notes. An issuer that marks its own bond at a high spread can offer better
terms. The note is still fairly priced if the spread is its real cost of funding.

After issue, the note lives on the desk's books like any derivative. The desk runs the \omterm{def:dv:greeks-and-the-hedging-pnl:greeks}{Greeks} of the option leg and funds or unfunds the
bond leg with treasury, and it quotes a secondary price, often as the issuer's buy-back. That price moves with the market and with the
issuer's spread: a widening of the spread lowers the value of every note outstanding, and an issuer that carries its notes at fair value
shows a gain on its own debt. Early redemptions, autocalls (chapter 18) and knock-ins change the funding the issuer holds, which the treasury has to plan for.

\section{Systematic-strategy indices and options on them}

Notes can also pay on a rules-based index rather than on a plain market index. Two such constructions are engineered to make options
cheaper.

\begin{definition}[Systematic strategy index]\label{def:dv:the-structured-products-business:ssi}
A \emph{systematic strategy index}\index{systematic strategy index} is an index whose level follows a published rule for trading one or
more underlyings, rebalanced mechanically, and on which notes and options are written.
\end{definition}

\begin{definition}[Volatility-target index]\label{def:dv:the-structured-products-business:vt}
A \emph{volatility-target index}\index{volatility-target index} holds an exposure to an underlying equal to a target volatility divided
by a recent estimate of the underlying's volatility, capped at a maximum leverage, with the rest in cash; its realised volatility stays
close to the target whatever the underlying does.
\end{definition}

\begin{definition}[Decrement index]\label{def:dv:the-structured-products-business:decrement}
A \emph{decrement index}\index{decrement index} is built on a total-return index by deducting a fixed amount, a percentage a year or a
number of points, in place of the dividends actually paid: its forward depends on that fixed deduction, not on uncertain dividend
forecasts.
\end{definition}

The volatility target makes the index's volatility nearly deterministic. An option on it is almost a Black--Scholes option at the target,
cheap when the target is below the market's volatility. The issuer's hedge then carries almost no volatility risk, and its \omterm{def:dv:greeks-and-the-hedging-pnl:greeks}{vega}
position no longer depends on the index's implied volatility. The \omterm{def:dv:the-structured-products-business:decrement}{decrement index} removes \omterm{def:dv:dividends-borrow-and-forwards:risk}{dividend risk} (chapter 5). The issuer knows
the forward exactly and can sell long-dated puts on it without holding dividend exposure. When the decrement exceeds the expected
dividends, the index drifts down relative to the market, and puts on it are dearer.

\begin{example}[Options on a 10\% volatility-target index]\label{ex:dv:the-structured-products-business:vt}
On chapter 10's \omterm{def:dv:stochastic-volatility:heston}{Heston model}, fitted to chapter 9's surface, build a 10\% \omterm{def:dv:the-structured-products-business:vt}{volatility-target index} with exponentially weighted volatility
estimates and leverage capped at 1.5, with zero rates. Over one year the raw index realises 19.8\% on average, with a standard deviation
of 8.2 points across paths. The target index realises 10.1\%, with a standard deviation of 0.6 point. A one-year at-the-money call costs
7.54 on the raw index and 4.07 on the target index, and the implied volatilities are 19.0\% and 10.2\%. The raw index's skew, from 21.6\% at
90 to 16.8\% at 110, nearly vanishes on the target index, 10.6\% to 9.9\% (\cref{fig:dv:the-structured-products-business:vt}).
\end{example}

\begin{omfigure}
\begin{tikzpicture}
\begin{groupplot}[group style={group size=2 by 1,horizontal sep=1.6cm},
  width=6.6cm,height=5cm,
  tick label style={font=\small},label style={font=\small},grid=major,grid style={gray!25},table/col sep=comma]
\nextgroupplot[xlabel={trading day},ylabel={level; exposure (\%)},xmin=0,xmax=252,
  title={\small one path},
  legend columns=1,legend cell align=left,legend style={font=\footnotesize,at={(0.5,-0.32)},anchor=north,draw=none,fill=none}]
\addplot[gray,thick] table[x=day,y=raw]{figdata/derivatives/19-the-structured-products-business/vt_path.csv};
\addplot[omThm,very thick] table[x=day,y=vt]{figdata/derivatives/19-the-structured-products-business/vt_path.csv};
\addplot[omDef,thick,dashed] table[x=day,y=exposure]{figdata/derivatives/19-the-structured-products-business/vt_path.csv};
\legend{index,10\% target index,exposure}
\nextgroupplot[xlabel={strike},ylabel={1-year implied volatility (\%)},xmin=87,xmax=113,ymin=5,ymax=25,xtick={90,100,110},
  title={\small calls on each},
  legend columns=1,legend cell align=left,legend style={font=\footnotesize,at={(0.5,-0.32)},anchor=north,draw=none,fill=none}]
\addplot[gray,very thick,mark=*,mark size=1.6pt] table[x=k,y=raw]{figdata/derivatives/19-the-structured-products-business/vt_iv.csv};
\addplot[omThm,very thick,mark=square*,mark size=1.6pt] table[x=k,y=vt]{figdata/derivatives/19-the-structured-products-business/vt_iv.csv};
\legend{index,10\% target index}
\end{groupplot}
\end{tikzpicture}

\omcaption{A 10\% \omterm{def:dv:the-structured-products-business:vt}{volatility-target index} on the Heston index of chapter 10. Left: one path of the index, the target index and
the exposure, which falls as realised volatility rises. Right: implied volatilities of one-year calls on each: the target index's are
close to 10\% at every strike. Data: the tutorial.}\label{fig:dv:the-structured-products-business:vt}
\end{omfigure}

\begin{example}[A decrement index against dividends]\label{ex:dv:the-structured-products-business:decrement}
Take a 3\% rate. A price index paying 3\% dividends has a five-year forward of 100. A 5\% \omterm{def:dv:the-structured-products-business:decrement}{decrement index} built on the same
total-return index has a forward of 90.5, since it loses 2\% a year relative to it
(\cref{fig:dv:the-structured-products-business:decrement}). A five-year at-the-money put at 20\% volatility costs 15.2 per 100 on the
price index and 19.0 on the \omterm{def:dv:the-structured-products-business:decrement}{decrement index}, 24\% more. That is what an investor who sells the put inside an \omterm{def:dv:autocallables:autocallable}{autocallable} receives in
exchange for bearing the gap between the decrement and the dividends actually paid.
\end{example}

\begin{omfigure}
\begin{tikzpicture}
\begin{axis}[width=10.5cm,height=4.6cm,
  xlabel={years},ylabel={forward (\% of today)},
  tick label style={font=\small},label style={font=\small},
  xmin=0,xmax=5,ymin=88,ymax=102,grid=major,grid style={gray!25},
  legend columns=2,legend cell align=left,legend style={font=\footnotesize,at={(0.5,-0.32)},anchor=north,draw=none,fill=none,/tikz/every even column/.append style={column sep=8pt}},
  table/col sep=comma]
\addplot[gray,very thick] table[x=t,y=price_fwd]{figdata/derivatives/19-the-structured-products-business/decrement.csv};
\addplot[omThm,very thick] table[x=t,y=dec_fwd]{figdata/derivatives/19-the-structured-products-business/decrement.csv};
\legend{price index (3\% dividends),5\% decrement index}
\end{axis}
\end{tikzpicture}

\omcaption{Forwards of a price index paying 3\% dividends and of a 5\% \omterm{def:dv:the-structured-products-business:decrement}{decrement index} on the same total-return index, with a 3\% rate.
The decrement fixes the forward, and the index drifts down against the market by the excess of the decrement over the dividends. Data:
the chapter's code.}\label{fig:dv:the-structured-products-business:decrement}
\end{omfigure}

Both constructions move a risk from the issuer to the investor. The target index gives the investor the timing risk of its rebalancing,
which buys after rallies and sells after falls, a short-gamma pattern. The decrement gives the investor the gap between the decrement and
the dividends. The payoff looks the same as on the plain index, and it is not. A structurer's first question about a strategy index is
therefore what it has taken out of the option price, and who now holds it.

\section{Tutorial: decomposing a note}\label{tut:dv:the-structured-products-business:tut}

\begin{tutorial}
\textbf{Goal.} Decompose a protected note into a bond at the issuer's rate and an option leg; solve for the participation, the cap, the
\omterm{def:dv:the-structured-products-business:rc}{reverse convertible}'s coupon and the funding spread that a target requires; build a volatility-target and a \omterm{def:dv:the-structured-products-business:decrement}{decrement index} and price
options on them. \textbf{End state:} the four figures and the numbers of the weekend problem.
\begin{tutsteps}
\item \textbf{The budget identity, both ways}:
\omcode{code/firm/termsheet/firm_termsheet.py}{57}{72}{Participation from the spread, and the spread from a participation.}\label{lst:dv:the-structured-products-business:budget}
\item \textbf{The \omterm{def:dv:the-structured-products-business:vt}{volatility-target index}}, with a lagged volatility estimate so the rule uses only past data:
\omcode{code/firm/termsheet/firm_termsheet.py}{86}{103}{Volatility estimate and target index.}\label{lst:dv:the-structured-products-business:vt}
\item \textbf{Run} \texttt{dv\_structured.named\_result()}, \texttt{capped\_note()}, \texttt{reverse\_convertible()},
\texttt{vol\_target\_options()}, \texttt{decrement\_example()} and \texttt{fig\_structured.py}.
\end{tutsteps}
\textbf{What to change next.} Lower the protection to 95\% and find the participation at a 1\% spread; replace the at-the-money call by the
average-price call of chapter 16; give the target index a 0.5\% annual fee and compare the call price.
\end{tutorial}

\section{Build: the term-sheet decomposer}\label{bld:dv:the-structured-products-business:termsheet}

\begin{build}
\bfield{Purpose} The miniature firm's term-sheet engine for its structured-products desk: notes as bond and option legs, the solvers that
turn a client's request into terms, and the strategy-index calculators.
\bfield{Interface} \texttt{Leg(kind, quantity, strike)}, \texttt{Note(maturity, legs)}; \texttt{zero\_coupon(t, r, spread)};
\texttt{value(note, r, spread, q, vol\_of\_k)}; \texttt{protected\_note(t, participation, protection, cap)};
\texttt{max\_participation(\ldots, margin, protection, cap)}; \texttt{spread\_for\_participation(target, \ldots)};
\texttt{reverse\_convertible\_coupon}; \texttt{ewma\_vol}, \texttt{vol\_target\_index(returns, target, rate, lam, max\_lev, dt, fee)};
\texttt{decrement\_index(total\_return, decrement, dt, points, base)}.
\bfield{Rules} The bond leg is always discounted at the issuer's funding curve; options at the desk's smile; the margin is an explicit input,
never a residual left in the option price; strategy indices use only information available at each rebalancing.
\bfield{Acceptance tests} \texttt{code/firm/termsheet/tests/}: the budget identity holds and inverts; a weaker issuer offers more
participation; a cap buys more participation; the \omterm{def:dv:the-structured-products-business:rc}{reverse convertible}'s coupon rises with the spread and the volatility; a target index on
constant volatility realises its target; a \omterm{def:dv:the-structured-products-business:decrement}{decrement index} on a deterministic total-return index has the expected forward, in percentage and
in points.
\bfield{Stretch} Notes on the \omterm{def:dv:autocallables:autocallable}{autocallable} pricer of chapter 18; secondary-market pricing with the issuer's current spread; the key information
document's performance scenarios.
\end{build}

\begin{omsources}
\item Regulation (EU) No 1286/2014 of the European Parliament and of the Council of 26 November 2014 on key information documents for
packaged retail and insurance-based investment products (PRIIPs), \emph{Official Journal of the European Union} L 352/1.
\item FINRA, Regulatory Notice 12-03, ``Heightened Supervision of Complex Products'' (January 2012).
\end{omsources}

\section{Exercises}

\begin{exercise}[$\star$]\label{exo:dv:the-structured-products-business:1}
Write a two-year 100\%-protected note with 60\% participation as a bond plus options. Which legs does the investor hold?
\end{exercise}

\begin{exercise}[$\star$]\label{exo:dv:the-structured-products-business:2}
With a 3\% rate, what is the two-year zero-coupon bond worth at spreads of 0 and 2.4\%? How much option budget does the spread add?
\end{exercise}

\begin{exercise}[$\star$]\label{exo:dv:the-structured-products-business:3}
Why is a note that offers 70\% participation not necessarily better value than one offering 50\%?
\end{exercise}

\begin{exercise}[$\star\star$]\label{exo:dv:the-structured-products-business:4}
Using the chapter's numbers (call 12.49, margin 1.5), what participation does a 95\%-protected note offer at a zero spread?
\end{exercise}

\begin{exercise}[$\star\star$]\label{exo:dv:the-structured-products-business:5}
Explain why options on a \omterm{def:dv:the-structured-products-business:vt}{volatility-target index} are cheap, and what the investor gives up.
\end{exercise}

\begin{exercise}[$\star\star$]\label{exo:dv:the-structured-products-business:6}
A 5\% \omterm{def:dv:the-structured-products-business:decrement}{decrement index} is built on a total-return index whose underlying pays 3\% dividends. By how much does it lag the price index each
year, and who bears the difference if dividends are cut to 1\%?
\end{exercise}

\begin{exercise}[$\star\star\star$]\label{exo:dv:the-structured-products-business:7}
\emph{Coding.} Price the \omterm{def:dv:the-structured-products-business:rc}{reverse convertible} of \cref{ex:dv:the-structured-products-business:rc} with strikes of 80\%, 90\% and 100\% and
report the coupons.
\end{exercise}

\begin{exercise}[$\star\star\star$]\label{exo:dv:the-structured-products-business:8}
\emph{Find the flaw.} ``Our notes offer the best participation on the market, so our option pricing must be the most competitive.''
\end{exercise}

\section{Problem: Seventy Percent Participation}

\begin{problem}[{Weekend problem --- what pays for the participation}]\label{pb:dv:the-structured-products-business:1}
A distributor asks for a two-year, 100\%-protected note on an index with 70\% participation. Rate 3\%, dividend yield 1.5\%, chapter 9's
surface, a \omterm{def:dv:the-structured-products-business:margin}{structuring margin} of 1.5.

\medskip\noindent\textbf{Part I --- The option leg.}
\begin{enumerate}
\item Give the two-year at-the-money volatility and the call price per 100.
\item What does 70\% participation cost?
\item Why is the call priced at the smile's at-the-money volatility and not at another strike's?
\item What would averaging the final level over the last six months do to the call's price?
\item What would a cap do?
\end{enumerate}

\medskip\noindent\textbf{Part II --- The bond leg.}
\begin{enumerate}[resume]
\item Give the zero-coupon bond's value at a zero spread.
\item Give the participation available at spreads of 0, 1\% and 2\%.
\item At what spread is 70\% possible? What does the bond cost there?
\item Why does the issuer's spread enter the note's price?
\item What does the investor bear in exchange for the higher participation?
\end{enumerate}

\medskip\noindent\textbf{Part III --- Alternatives.}
\begin{enumerate}[resume]
\item At a 1\% spread, what cap allows 100\% participation?
\item What does the same budget buy as a one-year \omterm{def:dv:the-structured-products-business:rc}{reverse convertible} struck at 90\%?
\item What participation would a 10\% \omterm{def:dv:the-structured-products-business:vt}{volatility-target index} allow, using the one-year implied volatility ratio as a guide?
\item Why might the distributor prefer a strategy index, and why might the investor not?
\item How does the key information document constrain the presentation?
\end{enumerate}

\medskip\noindent\textbf{Part IV --- Judgement.}
\begin{enumerate}[resume]
\item Is 70\% at 240 basis points a good deal for the client?
\item What would you tell the distributor about an issuer that offers 70\% when others offer 50\%?
\item How does a widening of the issuer's spread after issue affect the note's secondary price?
\item State the \emph{named result}: the issuer funding spread at which a two-year 100\%-protected note can offer 70\% participation after a
1.5\% margin.
\item In one sentence: what does a \omterm{def:dv:the-structured-products-business:protected}{capital-protected note}'s participation measure?
\end{enumerate}
\end{problem}

\section{Interview questions}

\begin{interviewq}[$\star$ \iqroles{trader, bank}]\label{iq:dv:the-structured-products-business:1}
Price a \omterm{def:dv:the-structured-products-business:protected}{capital-protected note}. What determines the \omterm{def:dv:the-structured-products-business:participation}{participation rate}?
\end{interviewq}

\begin{interviewq}[$\star$ \iqroles{bank}]\label{iq:dv:the-structured-products-business:2}
Why is a \omterm{def:dv:the-structured-products-business:note}{structured note} a funding instrument for the issuing bank?
\end{interviewq}

\begin{interviewq}[$\star\star$ \iqroles{trader}]\label{iq:dv:the-structured-products-business:3}
What is a \omterm{def:dv:the-structured-products-business:rc}{reverse convertible}, and what is the investor short?
\end{interviewq}

\begin{interviewq}[$\star\star$ \iqroles{researcher}]\label{iq:dv:the-structured-products-business:4}
Why are options on volatility-target indices cheaper than options on their underlying? What risk remains for the issuer?
\end{interviewq}

\begin{interviewq}[$\star\star$ \iqroles{risk}]\label{iq:dv:the-structured-products-business:5}
The issuer's credit spread widens by 100 basis points. What happens to its outstanding notes and to its accounts?
\end{interviewq}

\begin{interviewq}[$\star\star\star$ \iqroles{trader, bank}]\label{iq:dv:the-structured-products-business:6}
A distributor wants the highest possible coupon on a one-year note. Walk through the choices that raise it and what each costs the investor.
\end{interviewq}
